\section{\sys Design}

\subsection{Design Principles}

\sys follows three key design principles to address the limitations of existing GPU stacks:

\paragraph{Principle 1: Decouple stable mechanisms from dynamic programmable policies}
\noindent\textbf{Our Approach:} \sys explicitly adopts a principle of decoupling stable mechanisms from dynamically programmable policies. We strictly isolate the stable, minimal set of resource management primitives from implementation complexity into clearly defined and stable mechanism APIs exposed by the \sysdriver and runtime for device code. Policy logic is expressed separately as dynamically loadable eBPF programs, attaching to these stable mechanism hooks. This design trades off initial complexity in defining and implementing clear mechanism interfaces but gains substantial long-term flexibility and adaptability: new policies can be rapidly developed, experimented with, and deployed without invasive changes to underlying drivers or runtime components, greatly enhancing GPU datacenter agility and maintainability.

\paragraph{Principle 2: Unified CPU--GPU policy programmability interface and control plane}
\sys explicitly introduces a unified and symmetric CPU--GPU policy programmability paradigm. Rather than relying on disparate, incompatible policy frameworks on host and device sides, we mandate a single, consistent policy representation: all programmable logic across CPU and GPU boundaries is expressed using the same restricted subset of eBPF, with the same abstraction like virtual memory address and state management interface. Policies undergo unified static verification and execution models on both sides, enabling seamless migration, partitioning, and coordination of policy logic between host-side and device-side execution points. Although this unified paradigm intentionally restricts certain aspects of native GPU-side programmability (e.g., direct CUDA kernel manipulations), it provides substantial practical benefits in eliminating policy fragmentation, enabling straightforward cross-layer coordination, and greatly simplifying the development and deployment of adaptive GPU datacenter resource management policies.

\paragraph{Principle 3: Safety and efficiency through verified bounded execution and interfaces}
\noindent\textbf{Our Approach:} \sys explicitly prioritizes both safety and efficiency through a principle of verified bounded execution and interfaces. Rather than allowing arbitrary programmable logic, we restrict all policies—whether executed within GPU kernels or \sysdriver hooks—to a rigorously limited eBPF instructions and interfaces. A unified static verifier ensures that these policies are memory-safe, free from unbounded loops, guaranteed to terminate deterministically, and strictly bounded in resource usage and behavior. A lightweight runtime further enforces stringent per-invocation compute and memory limits, and employs GPU-specific optimizations to minimize overhead when executing device-side instrumentation. This design explicitly trades off some expressiveness in policy programming to ensure the reliability, determinism, and minimal latency overhead required in modern production GPU applications.

\subsection{High-Level Architecture}

\begin{figure}
\centering
\fbox{\parbox{0.8\columnwidth}{\centering [Figure: \sys Design]}}
\vspace{-.5ex}
\caption{\sys Design}
\label{fig:system-design}
\end{figure}

Figure~\ref{fig:system-design} shows the high-level architecture of \sys, clearly illustrating the CPU--GPU boundary, \sysdriver extension, and GPU kernel-level runtime integration points. \sys provides a unified eBPF runtime that manages verification and JIT compilation within the OS kernel. Source eBPF programs are compiled to standard bytecode using existing tools such as \texttt{clang}, \texttt{bpftool}, or \texttt{bpftrace}. Programs compiled with \texttt{clang}/\texttt{libbpf} write to eBPF maps in unified memory visible to both the OS kernel and GPU kernels. This bytecode is loaded into the OS kernel, where it undergoes standard verification for safety. For GPU execution, the kernel-mode runtime translates the verified bytecode into PTX and injects it into the target GPU context. On demand, eBPF bytecode is translated to PTX and injected at GPU hook points so device-side probes write into the same maps, yielding real-time, low-overhead cross CPU--GPU observability and control.

\subsection{Key Design Features}

In our design, the unified CPU--GPU policy programmability framework explicitly provides four novel technical capabilities:

\begin{itemize}[leftmargin=*,itemsep=0pt]
    \item \textbf{Unified CPU--GPU intermediate representation and execution model.}
    \sys adopts a single intermediate representation (a restricted subset of eBPF instructions and helpers) uniformly across CPU and GPU execution environments. This unified IR enables policy developers to use consistent programming abstractions, reasoning, and debugging techniques when implementing host-side and device-side logic, dramatically simplifying policy design and implementation complexity.

    \item \textbf{Unified cross-layer state management.}
    \sys provides unified, efficient cross-layer state management primitives accessible from both CPU- and GPU-side policy execution points. Policies can transparently maintain and exchange state (e.g., counters, histograms, policy flags) via common eBPF maps and shared data structures, without requiring explicit synchronization mechanisms or manual data transfer code between host and GPU device contexts. This significantly reduces the complexity associated with managing shared policy state across heterogeneous execution contexts.

    \item \textbf{Transparent cross-layer metadata and abstraction translation.}
    \sys transparently translates and propagates core execution metadata and abstractions---such as virtual memory addresses, kernel identifiers, timestamps, and execution contexts---across CPU and GPU environments. This translation provides policies with a consistent, unified view of crucial metadata and abstractions regardless of their execution location. As a result, developers can implement complex cross-layer policy logic without explicitly handling metadata translation, address translation, or synchronization across architectural boundaries.

    \item \textbf{Unified verification model and centralized policy control plane.}
    \sys implements a unified verification and safety model, ensuring that policy logic running on CPU and GPU sides is safe, bounded, and deterministic. A single verifier analyzes policy correctness, resource consumption, and execution guarantees before deployment. Additionally, \sys provides a centralized policy control plane that supports dynamic policy lifecycle management (loading, unloading, updating), adaptive parameter configuration, auxiliary computations, and runtime diagnostics. Together, this unified safety verification and centralized operational control enable developers to confidently deploy sophisticated policies onto latency-critical GPU and kernel paths with guaranteed safety and minimal overhead.
\end{itemize}

\subsection{\sysdevice Mechanism}\label{sec:\sys_device}

\noindent \sysdevice provides eBPF-style uprobe/uretprobe hooks for GPU functions at user level. The example below traces a tiled matrix multiply (\texttt{matMulTiled}), counting calls and measuring total time with standard eBPF maps:

\begin{minted}[highlightlines=false,fontsize=\small]{c}
SEC("\sys_device/_Z11matMulTiledPKfS0_Pf")
int \sys_device_matMulTiled(struct pt_regs *ctx) {
    u32 pid = bpf_get_current_pid_tgid() >> 32;
    u64 ts  = bpf_ktime_get_ns();
    bpf_map_update_elem(&start_ts, &pid, &ts, BPF_ANY);
    u64 one = 1, *cnt = bpf_map_lookup_elem(&call_count,
        &pid);  // increment call count
    if (cnt) *cnt += 1;
    else    bpf_map_update_elem(&call_count, &pid, &one,
        BPF_ANY);
}
SEC("\sys_device_ret/_Z11matMulTiledPKfS0_Pf")
int \sys_device_ret_matMulTiled(struct pt_regs *ctx) {
    u32 pid = bpf_get_current_pid_tgid() >> 32;
    u64 *tsp = bpf_map_lookup_elem(&start_ts, &pid);
    if (!tsp) return 0;
    u64 delta = bpf_ktime_get_ns() - *tsp;
    bpf_map_delete_elem(&start_ts, &pid);
    u64 *total = bpf_map_lookup_elem(&total_time_ns,
            &pid); // accumulate execution time
    if (total) {
        *total += delta;
        bpf_map_update_elem(&total_time_ns, &pid,
            total, BPF_EXIST);
    } else {
        bpf_map_update_elem(&total_time_ns, &pid,
            &delta, BPF_NOEXIST);
    }
}
\end{minted}

\noindent \sysdevice lifts the same semantics onto the device by JIT--translating probe handlers into PTX and injecting them at kernel launch time. Entry/exit trampolines record per-warp timestamps and update shared eBPF maps, preserving map semantics and enabling function-level profiling without application restarts.

\subsection{Extending \sysdriver (CPU-side)}

\sys safely extends \sysdriver on the CPU side by adding attach points within critical policy decision paths, including UVM memory management logic and GPU kernel scheduler logic. These extensions follow the struct\_ops-style callback approach used successfully in CPU-side eBPF systems, allowing developers to register verified eBPF programs at well-defined hook points without modifying core driver logic.

The eBPF verifier ensures that all attached programs are safe, bounded, and cannot corrupt driver state. CPU-side BPF maps provide efficient state coordination between the driver and eBPF policy logic, enabling adaptive decisions based on system-wide state while maintaining the safety guarantees of the eBPF verification framework.

\subsection{GPU-side eBPF Runtime Execution}

To enable fine-grained policy control within GPU kernels, \sys offloads verified eBPF programs to execute directly on the GPU device. This is achieved through a PTX-level JIT injection approach that translates verified eBPF bytecode into optimized PTX instructions and injects them at kernel launch.

The GPU runtime sandbox provides strong safety guarantees through bounded execution, careful register management, and memory access restrictions. Each injected program operates within a protected domain that prevents interference with application kernels while providing the flexibility needed for adaptive policy decisions.

\paragraph{GPU Trampoline}\label{sec:gpu_trampoline}
To instrument kernels transparently, \sys{} uses tiny PTX trampolines at entry/exit or hot sites. When reached, execution briefly diverts to our snippet and then resumes. Each trampoline:
\begin{itemize}[leftmargin=*,itemsep=0pt]
  \item \textbf{Save:} store minimal live state to a per‐warp scratch buffer.
  \item \textbf{Probe:} jump to the injected PTX (from eBPF) to read/update maps or record events.
  \item \textbf{Restore:} reload state and resume the original control flow.
\end{itemize}
Trampolines are patched at JIT time and confined to the protected domain (\S\ref{sec:isolation}); we observe sub--\(1\,\mu\)s overhead per invocation in practice.

\subsection{Unified CPU--GPU State Abstraction}

A key innovation in \sys is the unified state abstraction that spans CPU and GPU execution contexts. Shared eBPF maps are pinned in unified memory, providing a coherent view of policy-relevant state across both domains. This enables policies to make informed decisions based on global system state rather than being limited to local, domain-specific information.

\sys implements efficient synchronization and coherency management across the CPU-GPU boundary, with optimized paths for PCIe-based systems (requiring explicit memory barriers). This careful design ensures that cross-domain state updates are visible with minimal latency while preserving correctness guarantees.
